\documentclass[12pt, twoside]{article}
%\documentclass[12pt, twoside]{article}
\usepackage[letterpaper, margin=1in, headsep=0.2in]{geometry}
\setlength{\headheight}{0.6in}
%\usepackage[english]{babel}
\usepackage[utf8]{inputenc}
\usepackage{microtype}
\usepackage{amsmath}
\usepackage{amssymb}
%\usepackage{amsfonts}
\usepackage[nomessages]{fp} %\FPeval{\var-name}{2*sin(pi/6)}
\usepackage{siunitx} %units in math. eg 20\milli\meter
\usepackage{yhmath} % for arcs, overparenth command
\usepackage{tikz} %graphics
\usetikzlibrary{quotes, angles, arrows, arrows.meta}
\usepackage{pgfplots}
\usepgfplotslibrary{fillbetween}
\pgfplotsset{compat=1.16}
\usepackage{graphicx} %consider setting \graphicspath{{images/}}
\usepackage{parskip} %no paragraph indent
\usepackage{enumitem}
\usepackage{multicol}
\usepackage{venndiagram}

\usepackage{fancyhdr}
\pagestyle{fancy}
\fancyhf{}
\renewcommand{\headrulewidth}{0pt} % disable the underline of the header
\raggedbottom
\hfuzz=2mm %suppresses overfull box warnings

\usepackage{hyperref}

\fancyhead[LE]{\thepage}
\fancyhead[RO]{Markscheme}
\fancyhead[LO]{La Scuola d'Italia / Dr.\ Huson / IB Math 12 \\* 2 November 2026}

\begin{document}
\subsubsection*{3.1 Classwork: The Unit Circle and Trigonometric Graphs (no calculator) --- Markscheme}

\begin{enumerate}[itemsep=0.3cm]

%--- Problem 1: Label the unit circle from memory (16 points); exact values; Pythagorean identity, verify and use with quadrant [14 marks] ---
\item[1.] [Maximum mark: 14]

\textbf{(a)} Angles and coordinates $\left(\cos\theta,\ \sin\theta\right)$:

\begin{center}
\renewcommand{\arraystretch}{1.6}
\begin{tabular}{llll}
$0$: $(1,\ 0)$ &
$\frac{\pi}{2}$: $(0,\ 1)$ &
$\pi$: $(-1,\ 0)$ &
$\frac{3\pi}{2}$: $(0,\ -1)$ \\
$\frac{\pi}{6}$: $\left(\frac{\sqrt{3}}{2},\ \frac{1}{2}\right)$ &
$\frac{2\pi}{3}$: $\left(-\frac{1}{2},\ \frac{\sqrt{3}}{2}\right)$ &
$\frac{7\pi}{6}$: $\left(-\frac{\sqrt{3}}{2},\ -\frac{1}{2}\right)$ &
$\frac{5\pi}{3}$: $\left(\frac{1}{2},\ -\frac{\sqrt{3}}{2}\right)$ \\
$\frac{\pi}{4}$: $\left(\frac{\sqrt{2}}{2},\ \frac{\sqrt{2}}{2}\right)$ &
$\frac{3\pi}{4}$: $\left(-\frac{\sqrt{2}}{2},\ \frac{\sqrt{2}}{2}\right)$ &
$\frac{5\pi}{4}$: $\left(-\frac{\sqrt{2}}{2},\ -\frac{\sqrt{2}}{2}\right)$ &
$\frac{7\pi}{4}$: $\left(\frac{\sqrt{2}}{2},\ -\frac{\sqrt{2}}{2}\right)$ \\
$\frac{\pi}{3}$: $\left(\frac{1}{2},\ \frac{\sqrt{3}}{2}\right)$ &
$\frac{5\pi}{6}$: $\left(-\frac{\sqrt{3}}{2},\ \frac{1}{2}\right)$ &
$\frac{4\pi}{3}$: $\left(-\frac{1}{2},\ -\frac{\sqrt{3}}{2}\right)$ &
$\frac{11\pi}{6}$: $\left(\frac{\sqrt{3}}{2},\ -\frac{1}{2}\right)$ \\
\end{tabular}
\end{center}

\hfill \textbf{A1A1A1A1A1A1}

\emph{A1 for all 16 angles in radians at the correct points. A1 for the
four points on the axes. A1 for each quadrant's three non-axis points
(first, second, third, fourth), with both coordinates and both signs
correct. Accept $\frac{1}{\sqrt{2}}$ for $\frac{\sqrt{2}}{2}$. Angles in
degrees earn the first A0. Swapping the coordinates of the $\frac{\pi}{6}$
and $\frac{\pi}{3}$ points is the common error: A0 for that quadrant.}

\textbf{(b)} $\sin\frac{5\pi}{6} = \frac{1}{2}$ \hfill \textbf{A1}

$\cos\frac{4\pi}{3} = -\frac{1}{2}$ \hfill \textbf{A1}

$\tan\frac{7\pi}{4} = -1$ \hfill \textbf{A1}

\emph{Common errors, each A0: $\sin\frac{5\pi}{6} = -\frac{\sqrt{3}}{2}$
(the cosine), $\cos\frac{4\pi}{3} = \frac{1}{2}$ (sign), $\tan\frac{7\pi}{4}
= 1$ (sign).}

\textbf{(c)} $\sin^{2}\frac{5\pi}{6} + \cos^{2}\frac{5\pi}{6}
= \left(\frac{1}{2}\right)^{2} + \left(-\frac{\sqrt{3}}{2}\right)^{2}$
\hfill \textbf{M1}

$= \frac{1}{4} + \frac{3}{4}$ \hfill \textbf{A1}

$= 1$ \hfill \textbf{AG}

\emph{M1 for substituting both exact values. A1 for $\frac{1}{4}$ and
$\frac{3}{4}$ seen. Accept $\frac{\sqrt{3}}{2}$ without the negative
sign for the M1, since it is squared, but $-\frac{3}{4}$ for its square
earns A0.}

\textbf{(d)} $\cos^{2}\theta = 1 - \left(\frac{2}{3}\right)^{2}
= \frac{5}{9}$ \hfill \textbf{M1}

$\cos\theta = \pm\frac{\sqrt{5}}{3}$ \hfill \textbf{A1}

$\cos\theta = -\frac{\sqrt{5}}{3}$, since $\theta$ is in the second quadrant
\hfill \textbf{A1}

\emph{M1 for using the identity with $\sin\theta = \frac{2}{3}$. The final
A1 is for choosing the negative value; $+\frac{\sqrt{5}}{3}$ as the final
answer earns M1A1A0. Wrong methods: $1 - \frac{2}{3} = \frac{1}{3}$
(squaring forgotten) gives $\cos\theta = -\frac{1}{\sqrt{3}}$, earning
M0A0A0.}

\bigskip \hrule \bigskip

%--- Problem 2: Plot sine and cosine from unit-circle values; gradient of OP is tan(theta); line at pi/6 with intercept; line at 2pi/3 through a point [11 marks] ---
\item[2.] [Maximum mark: 11]

\textbf{(a)} Points $(0,\ 0)$, $\left(\frac{\pi}{6},\ 0.5\right)$,
$\left(\frac{\pi}{3},\ 0.87\right)$, $\left(\frac{\pi}{2},\ 1\right)$,
\ldots, $(\pi,\ 0)$, \ldots, $\left(\frac{3\pi}{2},\ -1\right)$, \ldots,
$(2\pi,\ 0)$ plotted \hfill \textbf{A1}

Smooth wave through the points, labelled $y = \sin x$ \hfill \textbf{A1}

\textbf{(b)} Points $(0,\ 1)$, $\left(\frac{\pi}{6},\ 0.87\right)$,
$\left(\frac{\pi}{3},\ 0.5\right)$, $\left(\frac{\pi}{2},\ 0\right)$,
\ldots, $(\pi,\ -1)$, \ldots, $(2\pi,\ 1)$ plotted \hfill \textbf{A1}

Smooth wave through the points, labelled $y = \cos x$ \hfill \textbf{A1}

\emph{In (a) and (b), the first A1 needs at least 11 of the 13 points
within half a grid square. The second A1 needs a smooth curve (no straight
segments, no corners at the maxima and minima) and a label. Curves
labelled the wrong way round lose the second A1 in both (a) and (b).}

\textbf{(c)} gradient of $OP = \dfrac{\sin\theta - 0}{\cos\theta - 0}$
\hfill \textbf{M1}

$= \dfrac{\sin\theta}{\cos\theta}$ \hfill \textbf{A1}

$= \tan\theta$ \hfill \textbf{AG}

\emph{M1 for rise over run using the coordinates of $O$ and $P$. A1 for
$\frac{\sin\theta}{\cos\theta}$ before the given answer. Accept a
right-triangle argument: opposite $\sin\theta$, adjacent $\cos\theta$.}

\textbf{(d)} $m = \tan\frac{\pi}{6} = \frac{1}{\sqrt{3}}$ \hfill \textbf{A1}

$y = \frac{1}{\sqrt{3}}x + 2$ \hfill \textbf{A1}

\emph{Accept $\frac{\sqrt{3}}{3}$ for the gradient. FT the equation from
the candidate's gradient. Wrong methods: $\tan\frac{\pi}{6} = \sqrt{3}$
(the value for $\frac{\pi}{3}$), or $m = \sin\frac{\pi}{6} = \frac{1}{2}$,
earn A0, then A1 FT for $y = mx + 2$.}

\textbf{(e)} $m = \tan\frac{2\pi}{3}$ \hfill \textbf{M1}

$= -\sqrt{3}$ \hfill \textbf{A1}

$y - 4 = -\sqrt{3}(x - 1)$ \quad (so $y = -\sqrt{3}x + 4 + \sqrt{3}$)
\hfill \textbf{A1}

\emph{Accept any correct form of the equation. FT the equation from the
candidate's gradient. Wrong methods: $m = +\sqrt{3}$ (sign) earns
M1A0A1 FT; swapping the coordinates of the point,
$y - 1 = -\sqrt{3}(x - 4)$, earns M1A1A0.}

\end{enumerate}
\end{document}
